\section{Background}\label{sec:background}
\subsection{Image Dataset Primitives}
In modern machine-learning frameworks, data-loading utilities separate sample access from batching and scheduling. In PyTorch, the \texttt{Dataset} class defines how to obtain a sample, while a \texttt{DataLoader} organizes iteration, batching, and parallel loading~\cite{pytorch}. Downloading is an earlier concern: the referenced bytes must first be acquired, associated with their metadata, and stored in a form that these utilities can consume. These stages should remain distinct so that a training pipeline can choose its own decoding, resizing, and augmentation policy.

Notably, the scale of image datasets has grown substantially in the modern era. CIFAR-10 contains 60,000 small images~\cite{cifar10}; ImageNet (1.2 million images) and ImageNet-21K (11 million images) support much larger classification tasks~\cite{imagenet,imagenet21k}. Finally, web-derived collections such as LAION-5B~\cite{laion5b} distribute over 5 billion images for training in state of the art models. At this scale, distributing a manifest is often more practical than redistributing every image. However, a URL is a reference to a resource, not a guarantee that the resource remains available or unchanged. Dataset construction must preserve the distinction between the requested population and the subset successfully acquired.

Two storage patterns are especially useful for machine learning training. A directory of individual images is convenient for inspection and small datasets, but very large file counts increase filesystem metadata overhead. Sharded tar archives provide a sequential representation that packages images with associated metadata and can support streaming training~\cite{tar,webdataset}. FLOW-DC separates output layout from concurrency control so that the same acquisition mechanism can serve different downstream conventions. In the evaluation here, original hosted bytes are preserved; image validity and archive integrity are checked separately.

\subsection{Parallel HTTP Data Transfer From Manifest}
A dataset manifest is a tabular description of remote objects and their associated metadata. At minimum, it identifies the URL for each requested item. Additional columns may contain labels, taxonomic identifiers, or source attribution. Manifest preparation determines the scientific population; acquisition determines which members can be delivered. If invalid rows were silently discarded, success would be measured against a smaller set than was requested, overstating completeness. FLOW-DC therefore keeps every original row identifiable through to the output and records an explicit disposition for each.

HTTP defines request and response semantics~\cite{rfc9110}. Its HTTP/1.1 mapping operates over a reliable transport~\cite{http}, commonly TCP, whose original specification~\cite{tcp} has been consolidated in a current standard~\cite{tcp-current}. A client issues a GET request, receives headers and a status, and reads the response body. A successful status alone does not establish that the saved file is usable: the body may be incomplete, invalid, or incorrectly associated with a manifest row. Transfer, persistence, and verification therefore contribute separately to reliable acquisition. Concurrent downloads overlap waiting time across requests. FLOW-DC uses Python asynchronous I/O: coroutines yield while waiting, and an event loop schedules other ready operations. Coroutines are a general mechanism for cooperative execution~\cite{async}; the HTTP implementation uses \texttt{aiohttp}~\cite{aiohttp}. This avoids dedicating a Python thread to every pending network operation.

The receiver exposes several useful signals. HTTP 429 explicitly indicates too many requests and can carry a \texttt{Retry-After} instruction~\cite{http-codes}. HTTP 503 can also include retry guidance~\cite{rfc9110}. Other server errors, connection failures, or rising delay may justify reducing demand, but they do not uniquely identify a rate limiter. In particular, HTTP 403 is not, on its own, proof that a client has been blocked for excessive concurrency. FLOW-DC combines conservative admission rules with outcome records rather than attempting to infer an undocumented quota from a single response.

Request concurrency, request-start rate, and byte throughput are different quantities. Large objects may keep a request active for a long time without a high start rate; small objects may do the opposite. Likewise, a client can maintain many open requests while delivering little verified output. We use concurrency to control admission and verified bytes per end-to-end second to assess useful delivery.

\subsection{BBR Inspiration}
FLOW-DC utilizes a congestion control algorithm known as PAARC (Policy-Aware Adaptive Request Control). PAARC draws inspiration mainly on Google's BBR (Bottleneck Bandwidth and Round-trip propagation time)~\cite{bbr} algorithm, which tries to run at full bandwidth before queues fill rather than waiting for packet loss to signal congestion. BBR models a transport path using an estimated bottleneck delivery bandwidth, $B$, and a propagation-delay reference, $R$~\cite{bbr}. Their product expresses an inflight-data scale:
\begin{equation}
\mathrm{BDP}=BR,\qquad r_{\mathrm{pace}}=g_pB,\qquad W=g_wBR,
\label{eq:bbr}
\end{equation}
where $g_p$ and $g_w$ are pacing and window gains. With bandwidth in bytes per second and delay in seconds, BDP and $W$ are byte quantities. The original BBR design separates startup, draining, bandwidth probing, and propagation-delay refresh, as summarized in Fig.~\ref{fig:bbr}.

\begin{figure*}[t]\centering
\includegraphics[width=.96\textwidth]{figures/bbr-phases.pdf}
\caption{Conceptual phases of the original BBR design~\cite{bbr}.}
\label{fig:bbr}\end{figure*}

PAARC adopts the broader principles of controlled exploration, a low-delay reference, and feedback-driven regulation. Its application-layer observations and actions are different. The measured signal is TTFB (Time To First Byte), the interval from dispatching request headers to receiving the first body byte on the final HTTP hop. It includes effects of network transit, server processing, and client scheduling; it is neither packet RTT nor a direct measurement of the server queue. The control variable is an integer request limit. PAARC does not derive its limit from a packet-level BDP. Delay-change signals also have established precedent in transport research, including TIMELY~\cite{timely}. Our contribution is the construction and integration of application-delay feedback with request admission, public-host safeguards, distributed execution, and verifiable dataset output. Table~\ref{tab:bbr} makes the intended correspondence explicit.

\begin{table*}[t]\centering\small
\caption{BBR inspiration and application-layer realization in PAARC}
\label{tab:bbr}
\begin{tabular}{p{.18\textwidth}p{.32\textwidth}p{.42\textwidth}}\toprule
Design element & Original BBR & PAARC and PAARC-Ratio\\\midrule
Controlled quantity & Packet pacing and inflight bytes & Per-host integer request-admission limit\\
Delay reference & Minimum packet round-trip estimate & Low-percentile application-delay baseline\\
Exploration & Probe transport delivery capacity & Grow a utilized limit during startup\\
Feedback & Delivery-rate and transport-path observations & Successful-response delay and explicit overload feedback\\
Refresh & Temporarily reduce inflight data & Reduce admission, drain requests, collect fresh samples\\
Regulation & Transport-model gains & Delay-ratio or persistent excess-delay-rate rule\\
Output objective & Efficient transport delivery & Verified original-row bytes per end-to-end second\\\bottomrule
\end{tabular}\end{table*}
